% Generated draft: requires mathematical and source review.
\documentclass[11pt]{article}
\usepackage[margin=25mm]{geometry}
\usepackage{fontspec}
\setmainfont{Latin Modern Roman}
\usepackage{amsmath,amssymb,mathtools,unicode-math}
\setmathfont{Latin Modern Math}
\usepackage{enumitem,needspace,xurl,hyperref}
\setlength{\parindent}{0pt}
\setlength{\parskip}{5pt}
\setlength{\emergencystretch}{3em}
\newcommand{\N}{\mathbb{N}}
\newcommand{\Z}{\mathbb{Z}}
\newcommand{\Q}{\mathbb{Q}}
\newcommand{\R}{\mathbb{R}}
\newcommand{\C}{\mathbb{C}}
\newcommand{\F}{\mathbb{F}}
\newcommand{\A}{\mathbb{A}}
\newcommand{\PP}{\mathbb{P}}
\newcommand{\E}{\mathbb{E}}
\newcommand{\eps}{\varepsilon}
\newcommand{\dd}{\,\mathrm{d}}
\begin{document}
\section*{UnsolvedMath ID: 701096}
\textbf{Source problem code:} catalogue entry
\subsection{Definitions and mathematical statement}
For a group \(G\) and a countable set \(X\), let \(\operatorname{Sym}(X)\) be the group of all bijections of \(X\). An action of \(G\) on \(X\) is a homomorphism \(\rho:G\to\operatorname{Sym}(X)\). It is faithful if \(\ker\rho=\{e\}\), and nontrivial if \(\rho(G)\ne\{\operatorname{id}_X\}\). Regarding a topological group as discrete means retaining its underlying abstract group and imposing no continuity requirement relative to its original topology on the action.

\textbf{Part (b), attributed to P. de la Harpe.} Is it true that for every Lie group \(G\), regarded as discrete, there exist a countable set \(X\) and an injective homomorphism
\[
\rho:G\longrightarrow\operatorname{Sym}(X)?
\]
The source further asks this for every \emph{separable continuous group}. It does not define that expression or specify the intended topological hypotheses, so the precise scope of this broader quantifier remains unresolved. The source also does not state its convention concerning countability hypotheses in the definition of a Lie group.

\textbf{Supplied qualifications for part (b).} The source's comment of 2021 reports an affirmative answer in the nilpotent case, citing N. Monod (2022). Its comment of 2025 reports a reduction to the case of simple Lie groups and, in particular, a solution of the solvable case, citing A. Conversano and N. Monod (2024). These are supplied literature assessments; their detailed hypotheses are not included in the record.

\textbf{Archived part (a), attributed to S. M. Ulam.} Let
\[
G=\operatorname{SO}(3)=\{A\in M_3(\R):A^{\mathsf T}A=I_3,\ \det A=1\},
\]
the usual compact rotation group of three-dimensional Euclidean space, regarded as discrete. Do there exist a countable set \(X\) and a homomorphism \(\rho:G\to\operatorname{Sym}(X)\) with \(\rho(G)\ne\{\operatorname{id}_X\}\)? The supplied archive reports an affirmative answer to this archived question, citing S. Thomas (1999) and another proof by Yu. L. Ershov and V. A. Churkin (2004). This is a supplied solved-part attribution, not an independently verified conclusion.
\subsection{Short English statement}
Part (b) asks whether every Lie group, considered as an abstract discrete group, has a faithful action on a countable set; the source proposes a broader class whose terminology needs clarification. Archived part (a) asks for a nontrivial countable action of \(\operatorname{SO}(3)\). The supplied record labels the combined problem partially solved and reports specific solved cases.
\subsection{Sources}
\begin{itemize}
\item Local UnsolvedMath registry, permanent ID \(701096\), exact source code \texttt{catalogue entry}, title \emph{Kourovka Notebook Problem 15.8}. The supplied record attributes the problem to P. de la Harpe, gives proposed year 2002, and attributes archived part (a) to S. M. Ulam.
\item \emph{The Kourovka Notebook (Unsolved Problems in Group Theory)}, 15th issue (2002), Problem 15.8; supplied background identifies an update in the 21st issue (2026), \texttt{arXiv:1401.0300v48}. The supplied collection description names E. I. Khukhro and V. D. Mazurov as editors.
\item Supplied Archive of Solved Problems reference for part (a): S. Thomas, \emph{J. Group Theory} \textbf{2} (1999), 401--434. The archive reports an affirmative answer.
\item Supplied archive reference for another proof of part (a): Yu. L. Ershov and V. A. Churkin, \emph{Dokl. Math.} \textbf{70}, no. 3 (2004), 896--898.
\item Supplied comment of 2021 concerning the nilpotent case: N. Monod, \emph{J. Group Theory} \textbf{25}, no. 5 (2022), 851--865.
\item Supplied comment of 2025 concerning reduction to simple Lie groups and settlement of the solvable case: A. Conversano and N. Monod, \emph{J. Algebra} \textbf{640} (2024), 106--116.
\item Supplied status review dated \texttt{2026-10-16}: canonical status \texttt{partially\_solved}, evidence label \texttt{published}; its notes place solved parts in the archive and remaining parts in the main section. This assessment is supplied dated curation, not independently verified.
\end{itemize}
\end{document}
